\section{Reading Comprehension Questions}

\noindent
{\em From Section~\ref{section:seq}}


\begin{requ}
If $\{x_n\}$ is defined by $x_n=3^n-2^n$, find $x_1$, $x_2$, $x_3$, $x_4$, and $x_5$.
\begin{requanswer}
$x_1=3-2=1$, $x_2=3^2-2^2=5$, $x_3=3^3-2^3=19$, $x_4=3^4-2^4=65$, and $x_5=3^5-2^5=211$.
\end{requanswer}
\end{requ}


\begin{requ}
What does it mean to {\em solve} a recurrence relation?
\begin{requanswer}
It means to find an explicit formula (or closed formula) for it.
In other words, a formula that is not recursively defined.
\end{requanswer}
\end{requ}

\begin{requ}
If $\{x_n\}$ is defined by $x_1=1$ and $x_n=2x_{n-1} + 3$, find $x_2$, $x_3$, $x_4$, and $x_5$.
\begin{requanswer}
$x_2=2 x_1+3 = 2\cdot 1+3=5$, $x_3=2 x_2+3 = 2\cdot 5+3=13$, $x_4=2 x_3+3=2\cdot 13+3=29$, 
and $x_5=2 x_4+3=2\cdot 29+3=61$.
\end{requanswer}
\end{requ}

\begin{requ}
Let $\{x_n\}$ be defined by $x_1=2$ and $x_n=x_{n-1}+3$.
\begin{lenum}
\item Find $x_2$, $x_3$, $x_4$, and $x_5$. 
\item Find a closed form for $x_n$.
\item Prove that your closed form is correct by following the technique from Example~\ref{exa:firstVerify}.
\end{lenum}
\begin{requanswer}
(a) $x_2=x_1+3 = 2+3=5$, $x_3=x_2+3 = 5+3=8$, $x_4=x_3+3=8+3=11$, and $x_5=x_4+3=11+3=14$.
(b) $x_n=3n-1$.
(c) Notice that if $x_n=3n-1$, then $x_{n-1}=3(n-1)-1=3n-4$. So
$x_{n-1}+3 = 3n-4+3=3n-1=x_n$, verifying that it works for the recurrence relation.
But we also need to show that it works for the base case: $x_1=2=3(1)-1$, so it also works for the base case and we
are done. 
\end{requanswer}
\end{requ}

\begin{requ}
Let $\{a_n\}$ be a sequence that gives the number of steps required to run some algorithm on an input of size $n$
(e.g. imagine the input is an array).
For instance, it takes $a_1$ steps to run the algorithm if the input is an array of size 1, 
$a_2$ steps if the input is an array of size 2, etc. 
Would you expect this sequence to be increasing, decreasing, or neither? Explain.
\begin{requanswer}
Increasing because for most algorithms, if you have more input it takes longer to even read the input, so it
would also take longer to run the algorithm. In rare cases it might be constant--for instance, an algorithm that
returns the fifth element of an array will take the same amount of time regardless of the size of the array.
\end{requanswer}
\end{requ}

\begin{requ}
Are geometric progressions always, sometime, or never monotonic? Explain.
Similarly, are they always, sometimes, or never increasing?
\begin{requanswer}
They are sometimes monotonic. If $r<0$ they will not be monotonic, but if $r>0$ they are monotonic.
They are also sometimes increasing. For instance, if $r>1$, it will be increasing, but if
$0<r<1$ it will be decreasing. If $r<0$, they are neither increasing nor decreasing since they
go back and forth between positive and negative.
\end{requanswer}
\end{requ}

\begin{requ}
Are arithmetic progressions always, sometime, or never monotonic? Explain.
Similarly, are they always, sometimes, or never increasing?
\begin{requanswer}
They are always monotonic. A given arithmetic progression is either always increasing or always decreasing,
depending on whether $d$ is positive or negative.
\end{requanswer}
\end{requ}

\begin{requ}
Give an example (not from the book) of each of the following.
\begin{lenum}
\item A geometric progression in closed form. 
\item An arithmetic progression in closed form.
\item A recurrence relation that defines a geometric progression.
\item A recurrence relation that defines an arithmetic progression.
\end{lenum}
\begin{requanswer}
Many answers are possible, but your answer should be very similar in form as the ones given.
(a) $a_n=3(2/3)^n$.
(b) $b_n=2+8n$.
(c) $c_n=(5/3)c_{n-1}$, $c_0=3$.
(d) $d_n=d_{n-1}+9/4$, $d_0=8$.
\end{requanswer}
\end{requ}

\noindent
{\em From Section~\ref{section:sumProd}}

\begin{requ}
Are $\dsum_{i=1}^{n} -x^i$ and $\dsum_{i=1}^{n} (-x)^i$ the same? Explain.
\begin{requanswer}
They are not because $-x^i=-(x^i)$ and $(-x)^i=(-1)^i(x)^i$. Depending on whether $x$ is positive or negative
and depending on whether $i$ is even or odd, these may have opposite signs.
\end{requanswer}
\end{requ}

\begin{requ}
Write $-1+3-9+27-81+243-729$ using a summation.
\begin{requanswer}
$\dsum_{i=0}^6 -(-3)^i$ or $\dsum_{i=0}^6 (-1)^{i+1}3^i$.
\end{requanswer}
\end{requ}

\begin{requ}
Compute $\dsum_{k=0}^{30} 5k-7$.
\begin{requanswer}
$\dsum_{k=0}^{30} 5k-7 = \dsum_{k=0}^{30} 5k - \dsum_{k=0}^{30} 7=5\dsum_{k=0}^{30} k - 7*31
=5\frac{30\cdot 31}{2} - 217 = 2108$
\end{requanswer}
\end{requ}

\begin{requ}
Compute $\dsum_{k=0}^{n} 2^k$. (Eventually you will hopefully have this one memorized.)
\begin{requanswer}
$\dsum_{k=0}^{n} 2^k = \frac{2^{n+1} -1}{2-1} = 2^{n+1}-1$ 
or 
$\dsum_{k=0}^{n} 2^k = \frac{1 - 2^{n+1}}{1 - 2} = \frac{1-2^{n+1}}{-1} = 2^{n+1}-1$. 
\end{requanswer}
\end{requ}

\begin{requ}
Is it ever the case that $\dsum_{i=0}^{n} x_i=\dsum_{i=1}^{n} x_i$? If so, when?
Give an example.
\begin{requanswer}
Sure. Whenever $x_0=0$.
\end{requanswer}
\end{requ}


\begin{requ}
Compute $\dsum_{k=1}^{23} \frac{11}{(-7)^k}$.  
Simplify your answer (although you don't need to compute the actual number). 
(I have thrown several subtle tricks at you on this one, but if you read carefully and apply what you know,
you should be able to do it!)
\begin{requanswer}
If you can get this one without mistakes, you are doing {\em really} well! If you don't quite get it, keep
trying to do it on your own. You will learn a lot in the process and it will be a good algebra refresher.

$\displaystyle
\dsum_{k=1}^{23} \frac{11}{(-7)^k} 
= 11\dsum_{k=1}^{23} \frac{1}{(-7)^k}
= 11\dsum_{k=1}^{23} \left(\frac{1}{-7}\right)^k
= 11\dsum_{k=1}^{23} \left(-\frac{1}{7}\right)^k
= 11\left(\dsum_{k=0}^{23} \left(-\frac{1}{7}\right)^k - \left(-\frac{1}{7}\right)^0\right)
$

$\displaystyle
= 11\left( \frac{1-\left(-\frac{1}{7}\right)^{24}}{1-\left(-\frac{1}{7}\right)} - 1\right)
= 11\left( \frac{1-\left(-1\right)^{24}\left(\frac{1}{7}\right)^{24}}{\frac{8}{7}} - 1\right)
= 11\left( \frac{7}{8}\left(1-\left(\frac{1}{7}\right)^{24}\right) - 1\right)
$

$\displaystyle
= 11\left(\frac{7}{8} -\frac{7}{8}\left(\frac{1}{7}\right)^{24} - 1\right)
= 11\left(-\frac{1}{8} -\frac{1}{8}\left(\frac{1}{7}\right)^{23}\right)
%= -\frac{11}{8} -\frac{11}{8}\left(\frac{1}{7}\right)^{23}
= -\frac{11}{8}\left(1+\left(\frac{1}{7}\right)^{23}\right)
%= \frac{77}{6}\left(1-\left(\frac{1}{7}\right)^{24}\right) - 11
%= \frac{11}{6}\left(\frac{1}{7}\right)^{23}+\frac{11}{6}
$.
\end{requanswer}
\end{requ}

\begin{requ}
Estimate $\cos (1)$ without using a calculator. 
\begin{requanswer}
According to Theorem~\ref{thm:maclaurin}, 
$\cos x = 1 - \dfrac{x^2}{2!} + \dfrac{x^4}{4!} - \cdots +(-1)^n\dfrac{x^{2n}}{(2n)!} +  \cdots$.
To approximate $\cos(1)$, we can just use several terms of the infinite sum.
So, 
$$\cos(1) \approx 1 - \dfrac{1^2}{2!} + \dfrac{1^4}{4!} -  \dfrac{1^6}{6!} 
=1-\dfrac{1}{2}+\dfrac{1}{24}-\dfrac{1}{720}
=\dfrac{720-360+30-1}{720} = \dfrac{389}{720}\approx 0.540277.$$
Note that $\cos(1)=0.540302305\cdots $, so our approximation is pretty good.
\end{requanswer}
\end{requ}


\noindent
{\em From Section~\ref{subsec:matDef}}


\begin{requ}
Can you add a $3\times 4$ matrix to a $4\times 3$ matrix? Explain.
\begin{requanswer} 
No. You can only add matrices that have the same 
dimensions (number of rows and columns).
\end{requanswer}
\end{requ}


\begin{requ}
Write out the $4\times 3$ matrix $A = [a_{ij}]$ where
$a_{ij} = 2^{i-1}3^{j-1}$.
\begin{requanswer} $A = \begin{bmatrix} 
1 & 3 & 9 \cr
2 & 6 & 18 \cr 
4 & 12 & 36 \cr 
8 & 24 & 72 \cr \end{bmatrix}.$  
\end{requanswer}
\end{requ}

\begin{requ}
Let 
$A=\begin{bmatrix} 
-3 & 0 & 1 \cr 
4 & -2 & 5 \cr\end{bmatrix}$
and 
$B=\begin{bmatrix} 
1 & 7 & -5 \cr 
7 & 1 & 8 \cr\end{bmatrix}$.
Compute $3A$ and $A-B$.
\begin{requanswer}
$3A=\begin{bmatrix} 
-9& 0 & 3 \cr 
12 & -6 & 15 \cr\end{bmatrix}$
$A-B=\begin{bmatrix} 
-4& -7 & 6 \cr 
-3 & -3 & -3 \cr\end{bmatrix}$.
\end{requanswer}
\end{requ}

\begin{requ}
Explicitly write out the entries of the matrices 
${\bf 0}_{2\times 5}$ and ${\bf I}_3$.
\begin{requanswer}
${\bf 0}_{3\times 4}=
\begin{bmatrix}0 & 0 & 0 & 0 &0 \cr 
0 & 0 & 0 & 0 &0\cr \end{bmatrix}$
and 
 ${\bf I}_3=
\begin{bmatrix}
1 & 0 & 0 \cr 
0 & 1 & 0 \cr 
0 & 0 & 1 \cr\end{bmatrix}$.
\end{requanswer}
\end{requ}

\noindent
{\em From Section~\ref{subsec:matMul}}


\begin{requ}
Can you multiply a $3\times 4$ matrix by a $4\times 7$ matrix?
What about multiplying a $6\times 3$ matrix by a $6\times 3$ matrix.
If you can, what are the dimensions of the result?
If you can't, explain why not.
\begin{requanswer} 
In the first case you can since the number of columns of
the first matrix is equal to the number or rows of the second one.
The result will have dimension $3\times 7$.
In the second case you cannot do the multiplication since $3\neq 6$.
That is, the number of columns of the first is not the same
as the number or rows of the second.
\end{requanswer}
\end{requ}

\begin{requ}
Determine the product $\displaystyle \begin{bmatrix} 1 & -1 \cr 1 &
1\cr\end{bmatrix}\begin{bmatrix}-2 & 1 \cr 0 & -1 \cr
\end{bmatrix}\begin{bmatrix} 1 & 1 \cr 1 & 2 \cr
\end{bmatrix}$.
\begin{requanswer}$\begin{bmatrix} 2 & 2 \cr 0 & -2 \cr
\end{bmatrix}$ 
\end{requanswer}
\end{requ}


%\begin{requ}
%Let $N = \begin{bmatrix}0 & -2  & -3 & -4 \cr 0 & 0 & -2 & -3  \cr
%0 & 0 & 0 & -2 \cr 0 & 0 & 0 & 0 \cr
%\end{bmatrix}$. Find $N^{2008}$.
%(Hint: Start by computing $N^2$, then $N^3$, $N^4$, and
%eventually you will be able to see where it is going.)
%\begin{requanswer}
%Notice $N^2= \begin{bmatrix}0 & 0  & 4 & 12
%\cr 0 & 0 & 0& 4  \cr  0 & 0 & 0 & 0 \cr 0 & 0 & 0 & 0 \cr
%\end{bmatrix}$,  $N^3= \begin{bmatrix}0 & 0  & 0 & -8 \cr 0 & 0 & 0&
%0  \cr  0 & 0 & 0 & 0 \cr 0 & 0 & 0 & 0 \cr
%\end{bmatrix}$ and $ N^4=\begin{bmatrix}0 & 0  & 0 & 0 \cr 0 & 0 & 0&
%0  \cr  0 & 0 & 0 & 0 \cr 0 & 0 & 0 & 0 \cr
%\end{bmatrix}$. Hence any power---from the fourth on---is the zero matrix.
%\end{requanswer}
%\end{requ}

\begin{requ}
Let $A$ be a $n\times n$ matrix such that $A^5={\bf 0}_n$.
What is $A^{19}$? 
\begin{requanswer} 
$A^{19}=A^5A^{14}={\bf 0}_n A^{14}={\bf 0}_n$.
\end{requanswer}
\end{requ}


\begin{requ}
Find all real numbers $x$ such that  
$$\begin{bmatrix} -4 & x \cr -x & 4
\end{bmatrix}^2 = \begin{bmatrix} -1 & 0 \cr 0 & -1
\end{bmatrix}.$$
\begin{requanswer} Observe that $$
\begin{bmatrix} -4 & x \cr -x & 4
\end{bmatrix}^2 =\begin{bmatrix} -4 & x \cr -x & 4
\end{bmatrix}\begin{bmatrix} -4 & x \cr -x & 4
\end{bmatrix} =  \begin{bmatrix} 16 - x^2 & 0 \cr 0 & 16 - x^2 \end{bmatrix}, $$ and so we must have
$16 - x^2 = -1$ which gives us $x = \pm\sqrt{17}$.
\end{requanswer}
\end{requ}

\begin{requ}
Prove or disprove: For all matrices $A,B\in\mats{n\times n}$,
$(A+B)(A-B) = A^2 - B^2.$
Hint: Play around with some simple $2\times 2$ matrices.
\begin{requanswer}
Disprove! Take for example $A = \begin{bmatrix} 0 & 0 \cr 1 & 1
\cr\end{bmatrix}$ and $B = \begin{bmatrix} 1 & 0 \cr 1 & 0
\cr\end{bmatrix}$. Then $$A^2 - B^2 = \begin{bmatrix} -1 & 0 \cr 0 &
1 \cr\end{bmatrix} \neq  \begin{bmatrix} -1 & 0 \cr -2 & 1
\cr\end{bmatrix}= (A+B)(A-B). $$
\end{requanswer}
\end{requ}


\noindent
{\em From Section~\ref{subsec:matT}}

\begin{requ}
Let $A, B \in \mats{2\times 2}$ be symmetric matrices.
Must their product $AB$ be symmetric? Prove or disprove!
\begin{requanswer}Disprove! This is not generally true. Take $A =
\dis{\begin{bmatrix} 1 & 1 \cr 1 & 2
\end{bmatrix}}$ and $B =
\dis{\begin{bmatrix} 3 & 0 \cr 0 & 1
\end{bmatrix}}$. Clearly $A^T = A$ and $B^T = B$. We have
$$AB  = \begin{bmatrix} 3 & 1 \cr 3 & 2\cr\end{bmatrix}$$but
$$(AB)^T  = \begin{bmatrix} 3 & 3 \cr 1 & 2\cr\end{bmatrix}.$$
\end{requanswer}
\end{requ}

\begin{requ}
Prove or disprove: If $A, B$ are square matrices of the same size,
then it is always true that $\tr{AB} = \tr{A}\tr{B}$.
\begin{requanswer} 
Disprove! Take $A = B = {\bf I}_n$ for any $n>1$. Then $\tr{AB} = n< n^2 = \tr{A}\tr{B}$.  \end{requanswer}
\end{requ}


\begin{requ}
Let $A$ be a square matrix. Prove that the matrix $AA^T$ is
symmetric.
\begin{requanswer} We have
$$(AA^T)^T = (A^T)^TA^T = AA^T.$$
\end{requanswer}
\end{requ}